\section{Autoencoder: aprender a normalidade comprimindo}
\label{sec:autoencoder}

\subsection{A ideia: comprimir e reconstruir}
\begin{definicao}[Autoencoder]
Uma rede que tenta copiar a própria entrada na saída, mas é \emph{obrigada a
passar por um gargalo}. Tem duas partes:
\begin{itemize}[noitemsep]
  \item \textbf{Encoder}: comprime a entrada em poucos números --- o
        \textbf{código} ou \emph{bottleneck} (gargalo).
  \item \textbf{Decoder}: tenta reconstruir a entrada original a partir desse
        código comprimido.
\end{itemize}
A perda é o quão diferente a reconstrução ficou da entrada (erro de reconstrução).
\end{definicao}

\begin{intuicao}
É como pedir a alguém para resumir um texto em uma frase e depois reescrever o
texto a partir só da frase. Para o resumo permitir uma boa reconstrução, ele
precisa capturar a \emph{essência} --- os padrões que se repetem. Detalhes
idiossincráticos se perdem no gargalo.
\end{intuicao}

\subsection{Por que isso detecta anomalias}
Aqui está o truque central do projeto. Treinamos o autoencoder \textbf{só com
dados normais}. Ele se torna especialista em reconstruir o normal --- e
\emph{apenas} o normal, porque foi só isso que viu.

\begin{itemize}[noitemsep]
  \item Entrada \textbf{normal} $\Rightarrow$ reconstrução boa $\Rightarrow$ erro
        \textbf{baixo}.
  \item Entrada \textbf{anômala} (padrão que ele nunca aprendeu) $\Rightarrow$
        reconstrução ruim $\Rightarrow$ erro \textbf{alto}.
\end{itemize}

\textbf{O erro de reconstrução vira o nosso medidor de anomalia.} Não precisamos
rotular nada: a própria incapacidade de reconstruir denuncia o anômalo. É isso
que torna o método \emph{não supervisionado}.

\begin{exemplo}[Fora de finanças]
Um autoencoder treinado só com fotos de gatos aprende a reconstruir gatos com
fidelidade. Mostre-lhe a foto de um carro: a reconstrução sai borrada e errada
(erro alto). Sem nunca ter recebido o rótulo ``isto não é gato'', ele acusa o
estranho pela qualidade da reconstrução.
\end{exemplo}

\subsection{A arquitetura do NeuraTrade}
Encadeamos LSTMs nas duas pontas, para respeitar a natureza sequencial dos dados:

\begin{lstlisting}[caption={LSTM-Autoencoder (esqueleto da arquitetura do projeto).}]
from tensorflow import keras
L = 30  # passos por janela
modelo = keras.Sequential([
    # ENCODER: comprime a janela em um vetor latente
    keras.layers.LSTM(64, input_shape=(L, 1)),
    keras.layers.Dropout(0.2),
    keras.layers.Dense(16, activation="relu"),   # gargalo (latent_dim=16)
    # DECODER: reconstroi a janela de 30 passos
    keras.layers.RepeatVector(L),
    keras.layers.LSTM(64, return_sequences=True),
    keras.layers.Dropout(0.2),
    keras.layers.TimeDistributed(keras.layers.Dense(1)),
])
modelo.compile(optimizer="adam", loss="mse")
\end{lstlisting}

\begin{definicao}[Dimensão latente (\emph{latent\_dim})]
O tamanho do gargalo. Pequeno demais: a rede não consegue nem reconstruir o
normal (perde sinal). Grande demais: o gargalo deixa de comprimir e a rede pode
aprender a copiar \emph{qualquer} coisa --- inclusive anomalias ---, perdendo
poder de detecção. É um equilíbrio.
\end{definicao}

\begin{naprat}
Gargalo de \textbf{16} (\texttt{latent\_dim=16}), LSTMs de 64 unidades,
\emph{dropout} $0{,}2$. Testamos $8/16/32$ e a perda de validação ficou
\emph{plana} entre eles --- ou seja, o resultado é pouco sensível nessa faixa ---,
o que confirmou 16 como escolha equilibrada (compressão de $\approx 4{:}1$ sobre
a janela de 30 passos). Treina-se \textbf{um modelo por ação}, para que cada um
aprenda a ``normalidade'' específica do seu ativo e setor.
\end{naprat}
